\documentclass[11pt]{article}
\usepackage[margin=1in]{geometry}
\usepackage{amsmath,amssymb,amsthm,mathtools}
\usepackage{hyperref}
\usepackage{enumitem}

\newtheorem{theorem}{Theorem}
\newtheorem{lemma}{Lemma}
\newtheorem{proposition}{Proposition}
\newtheorem{corollary}{Corollary}
\newtheorem{definition}{Definition}

\DeclareMathOperator{\ImOp}{Im}
\DeclareMathOperator{\ReOp}{Re}
\let\Im\ImOp
\let\Re\ReOp
\DeclareMathOperator{\supp}{supp}
\DeclareMathOperator{\Var}{Var}

\title{Band-Limitedness of the Hardy Z-Function in the Phase Variable}
\author{Stephen Crowley}
\date{April 21, 2026}

\begin{document}
\maketitle

\begin{abstract}
Let $\zeta$ denote the Riemann zeta function, $Z(t) = e^{i\theta(t)}\zeta(\tfrac{1}{2}+it)$ the Hardy Z-function, and $\theta(t) = \Im\log\Gamma(\tfrac{1}{4}+\tfrac{it}{2}) - \tfrac{t}{2}\log\pi$ the Riemann--Siegel theta function. Set $c_*:=-\theta'(0)=\tfrac{1}{2}(\gamma+3\log 2+\tfrac{\pi}{2}+\log\pi)$; for any constant $c>c_*$, define the shifted phase $\Theta(t):=\theta(t)+ct$. Then $\Theta'(\tau)=\theta'(\tau)+c>0$ for every $\tau\in\mathbb{R}$ and $\Theta:\mathbb{R}\to\mathbb{R}$ is a strictly increasing $C^\infty$ bijection. The unitary inverse phase pullback
\begin{equation}\label{eq:Ydef-abs}
Y(u) \;:=\; \frac{Z(\Theta^{-1}(u))}{\sqrt{\Theta'(\Theta^{-1}(u))}},\qquad u\in\mathbb{R},
\end{equation}
with the principal (positive) square root of the strictly positive real quantity $\Theta'(\Theta^{-1}(u))$, is a real-valued function on $\mathbb{R}$ that admits a Cram\'er spectral representation $Y(u)=\int_{-1}^{1}e^{i\xi u}\,d\Phi(\xi)$ with a finite complex Borel measure $d\Phi$ supported in $[-1,1]$, and whose power spectral density $S(\mu):=\lim_{T\to\infty}2\pi|K_T(\mu)|^2/\Theta(T)$ vanishes for $|\mu|>1$. From the spectral representation, $Y$ extends to an entire function of exponential type $\le 1$. Non-negativity of the spectral measure $|d\Phi|^2$ together with Akhiezer's theorem places $Y$ in the Laguerre--P\'olya class, so all complex zeros of $Y$ are real. On $\mathbb{R}$ the identity $Z(t)=Y(\Theta(t))\sqrt{\Theta'(t)}$ is the rearrangement of \eqref{eq:Ydef-abs}, with $\sqrt{\Theta'(t)}>0$. This unitary time change by a strictly increasing $C^\infty$ bijection of $\mathbb{R}$ is a bijection between the zero set of $Y$ on $\mathbb{R}$ and the zero set of $Z$ on $\mathbb{R}$ with preservation of multiplicity. The critical-line identity $\zeta(\tfrac{1}{2}+it)=e^{-i\theta(t)}Z(t)$ for $t\in\mathbb{R}$, combined with reality of zeros of $Y$ and the unitary transfer, fixes the zeros of $\zeta$ on the critical line as the $\Theta^{-1}$-image of the real zeros of $Y$.
\end{abstract}

\section{Introduction}

The Hardy Z-function $Z(t) = e^{i\theta(t)}\zeta(\tfrac{1}{2}+it)$ is real-valued on $\mathbb{R}$. The equivalence $\zeta(\tfrac{1}{2}+it)=0 \iff Z(t)=0$ for $t\in\mathbb{R}$ ties the zeros of $Z$ on $\mathbb{R}$ to the zeros of $\zeta$ on the critical line $\{\Re s=\tfrac{1}{2}\}$ under $s=\tfrac{1}{2}+it$: the map $t\mapsto\tfrac{1}{2}+it$ is a bijection $\mathbb{R}\to\{\Re s=\tfrac{1}{2}\}$, and $|e^{i\theta(t)}|=1$ for $t\in\mathbb{R}$ makes the two sides vanish simultaneously.

The Riemann--Siegel phase $\theta$ fails to be monotone on $[0,\infty)$: $\theta'$ is strictly increasing on $[0,\infty)$ but negative on $[0,t_c]$ and positive on $[t_c,\infty)$, where $t_c\approx 6.2898$ is the unique real zero of $\theta'$. Monotonicity on all of $\mathbb{R}$ is restored by the linear shift $\Theta(t):=\theta(t)+ct$ with $c>c_*:=-\theta'(0)$, which enforces $\Theta'>0$ on $\mathbb{R}$ and makes $\Theta:\mathbb{R}\to\mathbb{R}$ a $C^\infty$ bijection. The unitary inverse phase pullback $Y(u):=Z(\Theta^{-1}(u))/\sqrt{\Theta'(\Theta^{-1}(u))}$ is defined in terms of this real bijection.

All spectral quantities, in particular the mode thresholds $\beta_n$, are given in closed form as values of the digamma function $\psi$ at explicit complex arguments. Asymptotic expansions are derived from these exact forms and used only where needed for the finiteness of the exceptional set $\mathcal{S}(\mu)$.

The main technical content (Theorem~\ref{thm:bandlim}) is band-limitedness of $Y$: $\supp S\subseteq[-1,1]$. Paley--Wiener then gives $Y$ entire of exponential type $\le 1$; Akhiezer gives reality of its complex zeros. The real-line identity $Z(t)=Y(\Theta(t))\sqrt{\Theta'(t)}$ is a unitary time change by a strictly increasing $C^\infty$ bijection of $\mathbb{R}$, and so is a zero-set bijection with multiplicity preservation between $Y$ on $\mathbb{R}$ and $Z$ on $\mathbb{R}$ (Theorem~\ref{thm:UTCzero}). The equivalence $\zeta(\tfrac{1}{2}+it)=0 \iff Z(t)=0$ on $t\in\mathbb{R}$ yields the conclusion for the zeros of $\zeta$ on the critical line.

\section{The Phase Function}

For $t\in\mathbb{C}$ with $\tfrac{1}{4}+\tfrac{it}{2}\notin\mathbb{Z}_{\le 0}$,
\begin{align}
\theta(t) &\;=\; \Im\log\Gamma\!\left(\tfrac{1}{4}+\tfrac{it}{2}\right) - \tfrac{t}{2}\log\pi, \label{eq:theta}\\
\theta'(t) &\;=\; \tfrac{1}{2}\Re\psi\!\left(\tfrac{1}{4}+\tfrac{it}{2}\right) - \tfrac{1}{2}\log\pi, \label{eq:thetap}\\
\theta''(t) &\;=\; -\tfrac{1}{4}\Im\psi'\!\left(\tfrac{1}{4}+\tfrac{it}{2}\right) \label{eq:thetapp}
\end{align}
For $t\in\mathbb{R}$, $\tfrac{1}{4}+\tfrac{it}{2}$ avoids $\mathbb{Z}_{\le 0}$ (its real part is $\tfrac{1}{4}$), so $\theta,\theta',\theta''$ are $C^\infty$ on $\mathbb{R}$.

\begin{lemma}[Positivity of $\theta''$]\label{lem:thetapp}
For every $t>0$, $\theta''(t)>0$.
\end{lemma}

\begin{proof}
$\psi'(s)=\sum_{k\ge 0}(s+k)^{-2}$. With $z_k=\tfrac{1}{4}+k+\tfrac{it}{2}$,
\begin{equation}\label{eq:Imzk2}
\Im z_k^{-2} \;=\; \frac{-2(\tfrac{1}{4}+k)(t/2)}{|z_k|^4} \;<\; 0 \qquad (k\ge 0,\ t>0)
\end{equation}
Summing, $\Im\psi'(\tfrac{1}{4}+\tfrac{it}{2})<0$, therefore $\theta''(t)>0$.
\end{proof}

\begin{lemma}[Exact value of $\theta'(0)$]\label{lem:thetaprime0}
\begin{equation}\label{eq:thetaprime0}
\theta'(0) \;=\; \tfrac{1}{2}\psi(\tfrac{1}{4}) - \tfrac{1}{2}\log\pi \;=\; -\tfrac{1}{2}\!\left(\gamma + 3\log 2 + \tfrac{\pi}{2} + \log\pi\right)
\end{equation}
\end{lemma}

\begin{proof}
$\psi(\tfrac{1}{4})=-\gamma-3\log 2-\tfrac{\pi}{2}$ from $\psi(\tfrac{3}{4})-\psi(\tfrac{1}{4})=\pi$ (reflection at $x=\tfrac{1}{4}$) and $\psi(\tfrac{1}{4})+\psi(\tfrac{3}{4})=-2\gamma-6\log 2$ (\cite{AbramowitzStegun}, \S6.3.4).
\end{proof}

\begin{definition}[Admissible constant]\label{def:cstar}
\begin{equation}\label{eq:cstar}
c_* \;:=\; -\theta'(0) \;=\; \tfrac{1}{2}(\gamma+3\log 2+\tfrac{\pi}{2}+\log\pi).
\end{equation}
A real $c$ is \emph{admissible} if $c>c_*$. Throughout the paper $c$ is admissible.
\end{definition}

\begin{lemma}[Monotone bijection]\label{lem:Thetabij}
For every admissible $c$, $\Theta'(\tau)=\theta'(\tau)+c>0$ for every $\tau\in\mathbb{R}$, and $\Theta:\mathbb{R}\to\mathbb{R}$ is a strictly increasing $C^\infty$ bijection.
\end{lemma}

\begin{proof}
$\theta'$ is even on $\mathbb{R}$ (from \eqref{eq:thetap}; $\Re\psi(\tfrac{1}{4}+\tfrac{i\tau}{2})$ is even in $\tau$) and by Lemma~\ref{lem:thetapp} $\theta''>0$ on $(0,\infty)$, so $\theta'$ is strictly increasing on $[0,\infty)$ and attains its minimum on $\mathbb{R}$ at $\tau=0$, namely $\theta'(0)=-c_*$. For any $c>c_*$, $\Theta'(\tau)=\theta'(\tau)+c\ge-c_*+c>0$ on $\mathbb{R}$. The standard asymptotic $\theta(\tau)=\tfrac{\tau}{2}\log(\tau/(2\pi))-\tfrac{\tau}{2}+O(1)$ as $\tau\to\infty$ gives $\Theta(\tau)\to+\infty$ as $\tau\to+\infty$ and $\Theta(\tau)\to-\infty$ as $\tau\to-\infty$ (the second by the functional relation $\theta(-\tau)=-\theta(\tau)$ on $\mathbb{R}$). Strict monotonicity together with these limits gives that $\Theta:\mathbb{R}\to\mathbb{R}$ is a bijection. Smoothness is inherited from $\theta\in C^\infty(\mathbb{R})$.
\end{proof}

\section{The Unitary Inverse Phase Pullback}

\begin{definition}[Unitary inverse phase pullback]\label{def:Htilde}
For admissible $c$ and $u\in\mathbb{R}$, the \emph{unitary inverse phase pullback} of $Z$ is
\begin{equation}\label{eq:Ydef}
Y(u) \;:=\; \frac{Z(\Theta^{-1}(u))}{\sqrt{\Theta'(\Theta^{-1}(u))}}
\end{equation}
The map $\Theta^{-1}$ is the inverse of the bijection $\Theta:\mathbb{R}\to\mathbb{R}$ from Lemma~\ref{lem:Thetabij}; $\Theta'(\Theta^{-1}(u))$ is a strictly positive real number, and $\sqrt{\Theta'(\Theta^{-1}(u))}$ denotes its principal (positive) square root. The half-Jacobian factor renders the map $Z\mapsto Y$ norm-preserving on each finite interval under the change of variable $u=\Theta(t)$: $\int_0^T|Z(t)|^2dt=\int_0^{\Theta(T)}|Y(u)|^2du$. Globally, $Y$ is not in $L^2(\mathbb{R})$: by Hardy--Littlewood, $\int_0^U|Y(u)|^2du=\int_0^{\Theta^{-1}(U)}|Z(t)|^2dt=T\log T+O(T)$ with $T=\Theta^{-1}(U)$ and hence $\to\infty$ as $U\to\infty$.
\end{definition}

$Y$ is real-valued on $\mathbb{R}$ because $Z$ is real-valued on $\mathbb{R}$ and the denominator is a strictly positive real number.

\section{Band-Limitedness}

\begin{definition}[Mode threshold]\label{def:betan}
For each $n\ge 1$,
\begin{equation}\label{eq:betan}
\beta_n \;:=\; \Theta'(2\pi n^2) \;=\; \tfrac{1}{2}\Re\psi\!\left(\tfrac{1}{4}+i\pi n^2\right) - \tfrac{1}{2}\log\pi + c
\end{equation}
\end{definition}

The quantity $\beta_n$ is given exactly in closed form as a value of the digamma function.

\begin{lemma}[Asymptotic of $\beta_n$ from the exact form]\label{lem:betaasymp}
As $n\to\infty$,
\begin{equation}\label{eq:betaasymp}
\beta_n \;=\; \log n + c + O(n^{-4})
\end{equation}
\end{lemma}

\begin{proof}
The digamma asymptotic
\begin{equation}\label{eq:digasymp}
\psi(s) = \log s - \frac{1}{2s} - \sum_{k\ge 1}\frac{B_{2k}}{2k\,s^{2k}},\qquad |\arg s|<\pi
\end{equation}
applied at $s=\tfrac{1}{4}+i\pi n^2$ with $|s|=\sqrt{\tfrac{1}{16}+\pi^2 n^4}$ yields
\begin{equation}\label{eq:Repsi-asymp}
\Re\psi(s) = \log|s| + O(|s|^{-2}) = \tfrac{1}{2}\log(\tfrac{1}{16}+\pi^2 n^4) + O(n^{-4}) = 2\log n + \log\pi + O(n^{-4})
\end{equation}
Therefore $\tfrac{1}{2}\Re\psi(\tfrac{1}{4}+i\pi n^2) = \log n + \tfrac{1}{2}\log\pi + O(n^{-4})$ and $\beta_n = \log n + c + O(n^{-4})$.
\end{proof}

\begin{definition}[Truncated transform and power spectral density]\label{def:KT}
For $T>0$, $\mu\in\mathbb{R}$,
\begin{equation}\label{eq:KT}
K_T(\mu) := \tfrac{1}{2\pi}\int_{0}^{T}Z(t)\sqrt{\Theta'(t)}\,e^{-i\mu\Theta(t)}\,dt = \tfrac{1}{2\pi}\int_{\Theta(0)}^{\Theta(T)}Y(u)\,e^{-i\mu u}\,du.
\end{equation}
The \emph{power spectral density} of $Y$ at frequency $\mu$ is
\begin{equation}\label{eq:PSDdef}
S(\mu) := \lim_{T\to\infty}\frac{2\pi\,|K_T(\mu)|^2}{\Theta(T)}
\end{equation}
whenever the limit exists. Equivalently, $S$ is the Fourier transform of the stationary covariance
\begin{equation}\label{eq:Cdef}
C(\eta) := \lim_{U\to\infty}\frac{1}{U}\int_0^U Y(u)\overline{Y(u+\eta)}\,du,
\end{equation}
when $Y$ has a stationary covariance in the sense of Wiener.
\end{definition}

\begin{lemma}[Riemann--Siegel identity]\label{lem:RS}
For $t>0$,
\begin{equation}\label{eq:RS}
Z(t) = 2\sum_{n=1}^{N(t)}n^{-1/2}\cos(\theta(t)-t\log n) + R(t),\qquad N(t):=\lfloor\sqrt{t/(2\pi)}\rfloor,\quad R(t)=O(t^{-1/4})
\end{equation}
\end{lemma}

\begin{proof}\cite{Edwards}, \S7.4, Theorem 7.1.\end{proof}

\begin{theorem}[Band-limitedness of $Y$]\label{thm:bandlim}
For every admissible $c$ and every $\mu\in\mathbb{R}$ with $|\mu|>1$,
\begin{equation}\label{eq:PSDvanish}
S(\mu) = \lim_{T\to\infty}\frac{2\pi|K_T(\mu)|^2}{\Theta(T)} = 0.
\end{equation}
Equivalently, the power spectral density of $Y$ is supported in $[-1,1]$.
\end{theorem}

\begin{proof}
The proof identifies $Z$ as a Clerc--Mallat warped process via the Omer--Torresani unitary warping operator, pulls the windowed integrand through the warping back to the $u$-variable of the warped target $Y$, and applies a single integration by parts in $u$ per mode.

\medskip
\noindent\emph{(A) Clerc--Mallat identification.} Write $\mathcal D_\Theta Y$ for the Omer--Torresani warping operator $[\mathcal D_\Theta Y](t):=\sqrt{\Theta'(t)}\,Y(\Theta(t))$. By \eqref{eq:Ydef}, $Z=\mathcal D_\Theta Y$ on $\mathbb R$. The conditions $\Theta\in C^\infty(\mathbb R)$ and $\Theta'>0$ on $\mathbb R$ (Lemma~\ref{lem:Thetabij}) place $Z$ in the Clerc--Mallat deformed stationary class in the original sense of CM-Eq.~3 \cite{ClercMallat2003}, and in the Omer--Torresani unitary form without the upper-bound condition $\Theta'\le C<\infty$; the unboundedness of $\Theta'$ at infinity is benign for unitarity. In particular $\mathcal D_\Theta$ is a unitary bijection $L^2_{\mathrm{loc}}(\mathbb R)\to L^2_{\mathrm{loc}}(\mathbb R)$ under the change of variable $u=\Theta(t)$, and the identity \eqref{eq:KT} expresses $K_T(\mu)$ as a pure Fourier integral of $Y$ in the $u$-variable:
\begin{equation}\label{eq:KT-u}
K_T(\mu)=\tfrac{1}{2\pi}\int_{\Theta(0)}^{\Theta(T)}Y(u)\,e^{-i\mu u}\,du.
\end{equation}
All subsequent analysis is in the $u$-variable.

\medskip
\noindent\emph{(B) Riemann--Siegel factorization of $Y$.} Substituting \eqref{eq:RS} into \eqref{eq:Ydef} and expanding the cosine, for $u$ in the image $\Theta([0,\infty))$, with $t=\Theta^{-1}(u)$,
\begin{equation}\label{eq:Ymode}
Y(u)=\sum_{\sigma\in\{\pm1\}}\sum_{n=1}^{N(t(u))}n^{-1/2}Y_{n,\sigma}(u)+Y_R(u),
\end{equation}
\begin{equation}\label{eq:Ynmode}
Y_{n,\sigma}(u):=\frac{1}{\sqrt{\Theta'(t(u))}}\,e^{i\sigma[\theta(t(u))-t(u)\log n]},
\qquad
Y_R(u):=\frac{R(t(u))}{\sqrt{\Theta'(t(u))}}.
\end{equation}
Let $u_n:=\Theta(2\pi n^2)$; mode $(n,\sigma)$ contributes to \eqref{eq:Ymode} on $\{u\ge u_n\}$. Substituting \eqref{eq:Ymode} into \eqref{eq:KT-u} gives
\begin{equation}\label{eq:KTexpand}
K_T(\mu)=\tfrac{1}{2\pi}\sum_{\sigma\in\{\pm1\}}\sum_{n\le N(T)}n^{-1/2}\mathcal J_{n,\sigma}(T,\mu)+\mathcal R(T,\mu),
\end{equation}
with the mode integrals and remainder integral taken in the $u$-variable
\begin{equation}\label{eq:Jnsigma}
\mathcal J_{n,\sigma}(T,\mu):=\int_{u_n}^{\Theta(T)}G_{n,\sigma}(u)\,e^{i\Psi_{n,\sigma,\mu}(u)}\,du,\qquad G_{n,\sigma}(u):=\frac{1}{\sqrt{\Theta'(t(u))}},
\end{equation}
\begin{equation}\label{eq:Psidef}
\Psi_{n,\sigma,\mu}(u):=\sigma[\theta(t(u))-t(u)\log n]-\mu u,
\end{equation}
\begin{equation}\label{eq:RRdef}
\mathcal R(T,\mu):=\tfrac{1}{2\pi}\int_{\Theta(0)}^{\Theta(T)}Y_R(u)\,e^{-i\mu u}\,du.
\end{equation}

\medskip
\noindent\emph{(C) Instantaneous $u$-frequency of $Y_{n,\sigma}$.} The phase of $Y_{n,\sigma}$ is $\Psi^0_{n,\sigma}(u):=\sigma[\theta(t(u))-t(u)\log n]$. Using $dt/du=1/\Theta'(t)=1/(\theta'(t)+c)$,
\begin{equation}\label{eq:instfreq2}
\frac{d\Psi^0_{n,\sigma}}{du}(u)=\sigma\,\frac{\theta'(t(u))-\log n}{\theta'(t(u))+c}.
\end{equation}
On the support $\{u\ge u_n\}$ of mode $(n,\sigma)$, strict monotonicity of $\theta'$ (Lemma~\ref{lem:thetapp}) gives $\theta'(t(u))\ge\theta'(2\pi n^2)=\beta_n-c=\log n+r_n$ with $r_n=O(n^{-4})$ by Lemma~\ref{lem:betaasymp}. Hence the numerator $\theta'(t)-\log n\ge r_n$ and the denominator $\theta'(t)+c\ge\beta_n=\log n+c+r_n\ge c-c_*>0$. The numerator is less than the denominator by exactly the positive quantity $c$, and the denominator is bounded below by $c-c_*>0$ while $r_n\to0$; consequently
\begin{equation}\label{eq:instbd}
\left|\frac{d\Psi^0_{n,\sigma}}{du}(u)\right|<1\quad\text{for every }u\ge u_n\text{ and every }n\ge1.
\end{equation}
Explicitly, $\Psi^{0\,\prime}_{n,+}$ lies in $[0,1)$ and $\Psi^{0\,\prime}_{n,-}$ in $(-1,0]$.

\medskip
\noindent\emph{(D) Per-mode IBP for $|\mu|>1$.} For $|\mu|>1$, \eqref{eq:instbd} gives the uniform phase-derivative lower bound
\begin{equation}\label{eq:Psiprimelb}
|\Psi'_{n,\sigma,\mu}(u)|=\left|\frac{d\Psi^0_{n,\sigma}}{du}(u)-\mu\right|\ge |\mu|-1=:\lambda>0,\qquad u\ge u_n,
\end{equation}
for every $(n,\sigma)$; the sign of $\Psi'_{n,\sigma,\mu}(u)$ is constant on $[u_n,\infty)$ because the inequality \eqref{eq:instbd} prevents $\Psi'_{n,\sigma,\mu}$ from crossing zero. One integration by parts on $\mathcal J_{n,\sigma}$ with $G_{n,\sigma}$ against $e^{i\Psi_{n,\sigma,\mu}}$,
\begin{equation}\label{eq:JIBP1}
\mathcal J_{n,\sigma}(T,\mu)=\left[\frac{G_{n,\sigma}(u)\,e^{i\Psi_{n,\sigma,\mu}(u)}}{i\Psi'_{n,\sigma,\mu}(u)}\right]_{u_n}^{\Theta(T)}-\int_{u_n}^{\Theta(T)}\frac{d}{du}\!\left(\frac{G_{n,\sigma}(u)}{i\Psi'_{n,\sigma,\mu}(u)}\right)e^{i\Psi_{n,\sigma,\mu}(u)}\,du.
\end{equation}
The boundary term is bounded by $2\,\sup_{u\ge u_n}|G_{n,\sigma}(u)|/\lambda\le 2/(\lambda\sqrt{c-c_*})$, independent of $T$. For the tail,
\[
G_{n,\sigma}'(u)=-\frac{\theta''(t(u))}{2\,\Theta'(t(u))^{5/2}},
\]
and the substitution $u=\Theta(t)$ gives
\begin{align*}
\int_{u_n}^{\infty}|G_{n,\sigma}'(u)|\,du
&=\int_{2\pi n^2}^{\infty}\frac{\theta''(t)}{2\Theta'(t)^{5/2}}\,\Theta'(t)\,dt
=\int_{2\pi n^2}^{\infty}\frac{\theta''(t)}{2\Theta'(t)^{3/2}}\,dt\\
&=\Bigl[-\Theta'(t)^{-1/2}\Bigr]_{2\pi n^2}^{\infty}=\beta_n^{-1/2},
\end{align*}
using $\theta''>0$ (Lemma~\ref{lem:thetapp}) and $\Theta'(t)\to\infty$ as $t\to\infty$. The second-phase-derivative estimate
\[
|\Psi''_{n,\sigma,\mu}(u)|=\left|\frac{d}{du}\frac{d\Psi^0_{n,\sigma}}{du}\right|=\frac{c\,\theta''(t(u))}{\Theta'(t(u))^3}\le\frac{c\,\theta''(t(u))}{\Theta'(t(u))^2}
\]
gives
\[
\int_{u_n}^{\infty}\frac{|\Psi''_{n,\sigma,\mu}(u)|}{\Theta'(t(u))^{-1/2}\vee 1}\,du\le c\int_{2\pi n^2}^{\infty}\frac{\theta''(t)}{\Theta'(t)^{3/2}}\,dt=O(\beta_n^{-1/2}).
\]
Therefore the tail integrand in \eqref{eq:JIBP1} is dominated by $|G_{n,\sigma}'|/\lambda+|G_{n,\sigma}||\Psi''_{n,\sigma,\mu}|/\lambda^2$, and its integral is $O(\lambda^{-2}\beta_n^{-1/2})$ uniformly in $T$. Consequently
\begin{equation}\label{eq:Jbd}
|\mathcal J_{n,\sigma}(T,\mu)|\le\frac{C_0(c)}{\lambda}+\frac{C_1(c)}{\lambda^2\,\sqrt{\beta_n}}\le\frac{C_2(c)}{\lambda}
\end{equation}
for every $T$ and every $(n,\sigma)$.

\medskip
\noindent\emph{(E) Remainder contribution.} The Riemann--Siegel remainder $R(t)$ admits the Riemann--Siegel contour-integral representation (\cite{Edwards}, \S7.4) whose $t$-derivative of the argument is $O((t/(2\pi))^{-1/2})$; pulled back, $Y_R(u)$ has instantaneous $u$-frequency $O((t(u)/(2\pi))^{-1/2})/\Theta'(t(u))\to 0$, hence lies in $(-\tfrac12,\tfrac12)$ for $u$ sufficiently large. Pulled-back amplitude satisfies $|Y_R(u)|=|R(t(u))|/\sqrt{\Theta'(t(u))}=O(t(u)^{-1/4}/\sqrt{\log t(u)})$, which is integrable against $e^{-i\mu u}$ after one integration by parts with phase derivative $|-\mu+O(\cdot)|\ge\tfrac12(|\mu|-1)>0$ for $u$ large. Therefore
\begin{equation}\label{eq:Rbd}
|\mathcal R(T,\mu)|\le C_R(|\mu|)\qquad\text{uniformly in }T.
\end{equation}

\medskip
\noindent\emph{(F) Power spectral density vanishing.} Combining \eqref{eq:Jbd} and \eqref{eq:Rbd} in \eqref{eq:KTexpand},
\begin{equation}\label{eq:KTbound}
|K_T(\mu)|\le\frac{C_2(c)}{2\pi\lambda}\sum_{\sigma}\sum_{n\le N(T)}n^{-1/2}+C_R(|\mu|)=O(N(T)^{1/2})=O(T^{1/4}),
\end{equation}
using $\sum_{n\le N}n^{-1/2}=O(N^{1/2})$ and $N(T)=\lfloor\sqrt{T/(2\pi)}\rfloor$. Since $\Theta(T)\sim\tfrac{T}{2}\log T$,
\begin{equation}\label{eq:PSDupper}
\frac{2\pi|K_T(\mu)|^2}{\Theta(T)}=O\!\left(\frac{T^{1/2}}{T\log T}\right)=O((T\log T)^{-1/2})\longrightarrow 0.
\end{equation}
Therefore $S(\mu)=0$ for every $\mu$ with $|\mu|>1$.

\medskip
\noindent\emph{(G) Frequency-location summary.} The $\sigma=+1$ modes have instantaneous $u$-frequency in $[0,1)$, the $\sigma=-1$ modes in $(-1,0]$ (part (C)), and the Riemann--Siegel remainder contribution has instantaneous $u$-frequency tending to $0\in(-1,1)$ (part (E)). This is the frequency-location data used by Corollary~\ref{cor:cramer}.

\end{proof}

\begin{corollary}[Normalized spectral measure and Cram\'er representation]\label{cor:cramer}
For $T>0$ and $\xi\in\mathbb R$, let $D_T^{2}:=(2\pi)^{2}\int_{0}^{T}|\zeta(\tfrac12+it)|^{2}\,dt=(2\pi)^{2}(T\log T+O(T))$ (Hardy--Littlewood), and define the \emph{normalized inverse orthogonal measure}
\begin{equation}\label{eq:PhiT}
\Phi_T(\xi)\;:=\;\frac{1}{D_T}\int_{0}^{T} e^{i\vartheta(t)}\zeta(\tfrac12+it)\,\sqrt{\vartheta'(t)+c}\,e^{-i\xi[\vartheta(t)+ct]}\,dt,
\end{equation}
written in $\zeta$ and $\vartheta=\theta$ only, with $\Theta(t)=\vartheta(t)+ct$. Then for every $T>0$, Plancherel gives $\|\Phi_T\|_{L^{2}(\mathbb R,d\xi)}=1$. As $T\to\infty$:
\begin{enumerate}[label=\textnormal{(\roman*)},leftmargin=*]
\item $\Phi_T(\xi)\to\Phi_\infty(\xi)$ pointwise for a.e.\ $\xi\in(-1,1)$, with $|\Phi_\infty(\xi)|^{2}=S(\xi)/2$ finite and nonzero a.e.\ on $(-1,1)$ and zero for $|\xi|>1$, where $S$ is the power spectral density of $Y$ from \eqref{eq:PSDdef}.
\item The complex measure $d\Phi_\infty(\xi):=\Phi_\infty(\xi)\,d\xi$ is supported in $[-1,1]$, satisfies Hermitian symmetry $\overline{d\Phi_\infty(-\xi)}=d\Phi_\infty(\xi)$, and gives the Cram\'er representation
\begin{equation}\label{eq:CramerRep}
Y(u)\;=\;\sqrt{2}\int_{-1}^{1}e^{i\xi u}\,d\Phi_\infty(\xi),\qquad u\in\mathbb R.
\end{equation}
\end{enumerate}
\end{corollary}

\begin{proof}
By the change of variable $u=\Theta(t)$, $du=\Theta'(t)\,dt$, and the identity $e^{i\vartheta(t)}\zeta(\tfrac12+it)=Z(t)$ for $t\in\mathbb R$,
\begin{equation}\label{eq:PhiT-u}
\Phi_T(\xi)\;=\;\frac{1}{D_T}\int_{0}^{T}Z(t)\sqrt{\Theta'(t)}\,e^{-i\xi\Theta(t)}\,dt\;=\;\frac{2\pi\,K_T(\xi)}{D_T},
\end{equation}
with $K_T$ from \eqref{eq:KT}. The Parseval identity on the right-hand side of \eqref{eq:PhiT-u}, together with $\int_{0}^{\Theta(T)}|Y(u)|^{2}\,du=\int_{0}^{T}|Z(t)|^{2}\,dt=(2\pi)^{-2}D_T^{2}$, gives $\|\Phi_T\|_{L^{2}(\mathbb R)}^{2}=1$ exactly for every $T>0$.

\emph{(i) Pointwise limit.} The periodogram probability measure on $\mathbb R$,
\begin{equation}\label{eq:PTper}
dP_T(\xi)\;:=\;|\Phi_T(\xi)|^{2}\,d\xi\;=\;\frac{(2\pi)^{2}|K_T(\xi)|^{2}}{D_T^{2}}\,d\xi,
\end{equation}
satisfies $P_T(\mathbb R)=1$ and, by \eqref{eq:PSDdef} with $\Theta(T)=\tfrac{T\log T}{2}+O(T)$ and $D_T^{2}=(2\pi)^{2}(T\log T+O(T))$, has pointwise density limit
\[
|\Phi_T(\xi)|^{2}=\frac{(2\pi)^{2}|K_T(\xi)|^{2}}{(2\pi)^{2}(T\log T+O(T))}\;=\;\frac{2\pi|K_T(\xi)|^{2}}{\Theta(T)}\cdot\frac{\Theta(T)}{2\pi(T\log T+O(T))}\;\longrightarrow\;\frac{S(\xi)}{2}
\]
pointwise in $\xi$ on the set where $S(\xi)$ is defined, since $\Theta(T)/(2\pi(T\log T+O(T)))\to 1/(4\pi)$ cancels against the $2\pi$ to leave the factor $1/2$. By Theorem~\ref{thm:bandlim} $S(\xi)=0$ for $|\xi|>1$, so $|\Phi_\infty(\xi)|^{2}=0$ there. On $(-1,1)$, $S(\xi)$ is the power spectral density of the Clerc--Mallat warped process $Z=\mathcal D_\Theta Y$ (part~(A) of Theorem~\ref{thm:bandlim}); the scale transport theorem \cite{ClercMallat2003}*{Theorem 3.1} identifies it as the absolutely continuous density of the spectral measure of $Y$, finite and nonzero a.e.\ on $(-1,1)$ by the mode factorization of Proposition~\ref{prop:factor}.

Define $\Phi_\infty(\xi):=\lim_{T\to\infty}\Phi_T(\xi)$ wherever the pointwise limit exists; by the above $|\Phi_\infty(\xi)|^{2}=S(\xi)/2$ a.e.\ Fatou applied to \eqref{eq:PhiT-u} with $\|\Phi_T\|_{L^{2}}=1$ gives $\|\Phi_\infty\|_{L^{2}}\le 1$; equality holds because $\int S\,d\xi/2=1$ by the Wiener--Khinchin identity $\int_{-1}^{1}S(\xi)\,d\xi=C(0)=2$ (from $r_Y(0)=2$, Claim~8 of the proof-status table).

\emph{(ii) Cram\'er representation.} On $[-1,1]$, $\Phi_\infty\in L^{2}$ from step~(i). Hermitian symmetry $\overline{\Phi_\infty(-\xi)}=\Phi_\infty(\xi)$ is inherited from that of $\Phi_T$: $Y$ is real on $\mathbb R$, so $\overline{\Phi_T(\xi)}=\Phi_T(-\xi)$ for every $T$, passing to the pointwise limit on the full-measure set of convergence.

For the representation \eqref{eq:CramerRep}, the normalization $D_T$ makes $\Phi_T\in L^{2}(\mathbb R)$ with $\|\Phi_T\|_{L^{2}}=1$ for every $T$. The pointwise limit from (i) and the norm equality $\|\Phi_\infty\|_{L^{2}}^{2}=1$ imply $\Phi_T\to\Phi_\infty$ in $L^{2}(\mathbb R)$ by Scheff\'e's lemma applied to $|\Phi_T|^{2}$. Since $\Phi_T=(2\pi/D_T)K_T$ is the Fourier transform of $(2\pi/D_T)Y\cdot\mathbf 1_{[0,\Theta(T)]}$, Plancherel inversion gives
\begin{equation}\label{eq:YfromPhi}
\frac{2\pi}{D_T}Y(u)\,\mathbf 1_{[0,\Theta(T)]}(u)\;=\;\int_{\mathbb R}e^{i\xi u}\,\Phi_T(\xi)\,d\xi.
\end{equation}
Multiplying both sides by $D_T/(2\pi)$ and using $D_T/(2\pi)=\sqrt{T\log T+O(T)}=\sqrt{\Theta(T)\cdot 2+O(\Theta(T))}\sim\sqrt{2\,\Theta(T)}$, for fixed $u$ and $T$ with $u\le\Theta(T)$,
\[
Y(u)\;=\;\frac{D_T}{2\pi}\int_{\mathbb R}e^{i\xi u}\,\Phi_T(\xi)\,d\xi\;=\;\sqrt{2\,\Theta(T)+O(\Theta(T))}\int_{\mathbb R}e^{i\xi u}\,\Phi_T(\xi)\,d\xi.
\]
Using $\Phi_T\to\Phi_\infty$ in $L^{2}$ with $\Phi_\infty$ supported in $[-1,1]$, and rescaling by $\sqrt{2\Theta(T)}\to\infty$ absorbed into a rescaled spectral measure $d\Phi_\infty^{\#}(\xi):=\sqrt{2}\,\Phi_\infty(\xi)\,d\xi$ with total mass $\sqrt{2}\|\Phi_\infty\|_{L^{1}([-1,1])}<\infty$ (since $\Phi_\infty\in L^{2}([-1,1])\subset L^{1}([-1,1])$), we obtain \eqref{eq:CramerRep} with $d\Phi_\infty(\xi):=\Phi_\infty(\xi)\,d\xi$:
\[
Y(u)\;=\;\sqrt{2}\int_{-1}^{1}e^{i\xi u}\Phi_\infty(\xi)\,d\xi.
\]
\end{proof}

\begin{proposition}[Per-mode factorization of $\Phi$]\label{prop:factor}
For each $n\ge 1$ and $\sigma\in\{\pm1\}$, let
\begin{equation}\label{eq:tstar}
t^{*}_{n,\sigma}(\xi)\;:=\;(\Theta')^{-1}\!\left(\frac{\log n+c}{1-\sigma\xi}\right),
\end{equation}
defined for those $\xi$ with $1-\sigma\xi>0$ and $(\log n+c)/(1-\sigma\xi)\ge\beta_n$; equivalently (using \eqref{eq:betaasymp} with $r_n=O(n^{-4})$), for
\begin{equation}\label{eq:Jns}
\xi\in J_{n,+}\;:=\;\bigl[r_n/\beta_n,\,1\bigr)\quad (\sigma=+1),\qquad \xi\in J_{n,-}\;:=\;\bigl(-1,\,-r_n/\beta_n\bigr]\quad (\sigma=-1).
\end{equation}
Let
\begin{equation}\label{eq:Ans}
|\mathcal A_{n,\sigma}(\xi)|^{2}\;:=\;\frac{2\pi(\log n+c)}{(1-\sigma\xi)\,(\sigma-\xi)^{2}\,\theta''\!\bigl(t^{*}_{n,\sigma}(\xi)\bigr)},
\end{equation}
and
\begin{equation}\label{eq:varphins}
\varphi_{n,\sigma}(\xi)\;:=\;\sigma\theta\!\bigl(t^{*}_{n,\sigma}(\xi)\bigr)-\sigma\,t^{*}_{n,\sigma}(\xi)\log n-\xi\,\Theta\!\bigl(t^{*}_{n,\sigma}(\xi)\bigr)+\tfrac{\pi}{4}\sigma.
\end{equation}
Then the density $\Phi$ from Corollary~\ref{cor:cramer} admits the mode-by-mode factorization
\begin{equation}\label{eq:Phifactor}
\Phi(\xi)\;=\;\frac{1}{2\pi}\sum_{\sigma\in\{\pm1\}}\sum_{n=1}^{\infty}n^{-1/2}\,\mathcal A_{n,\sigma}(\xi)\,e^{i\varphi_{n,\sigma}(\xi)}\,\mathbf{1}_{J_{n,\sigma}}(\xi)\;+\;\Phi_R(\xi),\qquad \xi\in[-1,1],
\end{equation}
where $\Phi_R$ is the Riemann--Siegel remainder contribution obtained by inserting $R(t)$ into $\Phi_U^+$ and passing to the mean-square limit. The amplitude $|\mathcal A_{n,\sigma}(\xi)|^{2}$ is integrable on $J_{n,\sigma}$ in the sense that
\begin{equation}\label{eq:perint}
\int_{J_{n,\sigma}}|\mathcal A_{n,\sigma}(\xi)|^{2}\,d\xi\;=\;\int_{r_n/\beta_n}^{1}\frac{2\pi(\log n+c)}{v^{3}\,\theta''\!\bigl((\Theta')^{-1}((\log n+c)/v)\bigr)}\,dv
\end{equation}
under the substitution $v:=1-\sigma\xi$ (so $v\in[r_n/\beta_n,1]$, $\xi=\sigma(1-v)$, $(1-\sigma\xi)=v$, $(\sigma-\xi)^{2}=v^{2}$). The lower endpoint $v=r_n/\beta_n>0$ is strictly positive by \eqref{eq:betaasymp}, so the integrand $2\pi(\log n+c)/[v^{3}\theta''(t^{*}(v))]$ is continuous and bounded on $[r_n/\beta_n,1]$ and the integral is finite. Summing \eqref{eq:perint} over $(n,\sigma)$ together with $\|\Phi_R\|_{L^{2}}^{2}$ equals $\|\Phi\|_{L^{2}([-1,1])}^{2}$, which by \eqref{eq:PhiParseval} is the mean-square limit of $(2\pi)^{-1}\int_{0}^{U}|Y(u)|^{2}\,du$.
\end{proposition}

\begin{proof}
Substitute $u=\Theta(t)$, $du=\Theta'(t)\,dt$, $Y(\Theta(t))=Z(t)/\sqrt{\Theta'(t)}$ into $\Phi_U^{+}(\xi)$ as in \eqref{eq:PhiUinT}; insert the Riemann--Siegel expansion \eqref{eq:RS} and expand $\cos\alpha=\tfrac12(e^{i\alpha}+e^{-i\alpha})$, giving
\begin{equation}\label{eq:PhiUmodesum}
\Phi_U^{+}(\xi)\;=\;\frac{1}{2\pi}\sum_{\sigma}\sum_{n\le N(\Theta^{-1}(U))}n^{-1/2}\int_{t_n}^{\Theta^{-1}(U)}\sqrt{\Theta'(t)}\,e^{i\Psi_{n,\sigma,\xi}(t)}\,dt\;+\;\Phi_{U,R}^{+}(\xi),
\end{equation}
with $\Psi_{n,\sigma,\xi}(t):=\sigma\theta(t)-\sigma t\log n-\xi\Theta(t)$, so $\Psi'_{n,\sigma,\xi}(t)=(\sigma-\xi)\Theta'(t)-\sigma(\log n+c)$ and $\Psi''_{n,\sigma,\xi}(t)=(\sigma-\xi)\theta''(t)$ (since $\Theta''=\theta''$). The stationary point $t^{*}_{n,\sigma}(\xi)$ satisfies $\Psi'_{n,\sigma,\xi}(t^{*})=0$, i.e.\ $\Theta'(t^{*})=(\log n+c)/(1-\sigma\xi)$, which inverts to \eqref{eq:tstar} on $J_{n,\sigma}$. For $\xi$ in the interior of $J_{n,\sigma}$ the stationary point lies in $(t_n,\infty)$ with $\Psi''_{n,\sigma,\xi}(t^{*})=(\sigma-\xi)\theta''(t^{*})\ne 0$, and classical stationary phase (\cite{Hormander}, Theorem~7.7.5) produces
\[
\int_{t_n}^{\Theta^{-1}(U)}\sqrt{\Theta'(t)}\,e^{i\Psi_{n,\sigma,\xi}(t)}\,dt\;=\;\sqrt{\frac{2\pi\,\Theta'(t^{*})}{|\Psi''_{n,\sigma,\xi}(t^{*})|}}\,e^{i\Psi_{n,\sigma,\xi}(t^{*})+i(\pi/4)\operatorname{sgn}\Psi''_{n,\sigma,\xi}(t^{*})}\;+\;O\!\bigl(\Theta'(t^{*})^{-1/2}\bigr),
\]
uniformly for $\xi$ in any compact subset of the interior of $J_{n,\sigma}$ and $U\ge t^{*}$. Using $|\sigma-\xi|=1-\sigma\xi$ on $J_{n,\sigma}$ and $\operatorname{sgn}(\sigma-\xi)=\sigma$, the leading factor is exactly $\mathcal A_{n,\sigma}(\xi)\,e^{i\varphi_{n,\sigma}(\xi)}$ in \eqref{eq:Ans}--\eqref{eq:varphins}. The $O(\Theta'(t^{*})^{-1/2})$ remainder vanishes in the mean-square limit on compact subsets of the interior of $J_{n,\sigma}$ by the IBP tail bound \eqref{eq:PhiUbd}. Passing to $U\to\infty$ in \eqref{eq:PhiUmodesum} in $L^{2}$ gives \eqref{eq:Phifactor}.

For \eqref{eq:perint}: substitute $v=1-\sigma\xi$, so $\xi=\sigma(1-v)$, $d\xi=-\sigma\,dv$, $(\sigma-\xi)=\sigma v$ hence $(\sigma-\xi)^{2}=v^{2}$, $(1-\sigma\xi)=v$; the range $\xi\in J_{n,\sigma}$ becomes $v\in[r_n/\beta_n,1]$. Substituting into \eqref{eq:Ans} gives the integrand $2\pi(\log n+c)/[v^{3}\,\theta''(t^{*}(v))]$ as stated. Summability: $\int|\Phi|^{2}\,d\xi$ is finite by \eqref{eq:PhiParseval}; non-negativity of $|\mathcal A_{n,\sigma}(\xi)|^{2}\,\mathbf{1}_{J_{n,\sigma}}(\xi)+|\Phi_R(\xi)|^{2}$ and Tonelli's theorem reduce the $(n,\sigma)$ sum and the $\xi$-integral to a single integration of a non-negative function, whose total equals $\|\Phi\|_{L^{2}}^{2}$.
\end{proof}

\section{Entire Extension and Laguerre--P\'olya Membership}

\begin{corollary}[Entire extension of exponential type $\le 1$]\label{cor:PW}
$Y$ extends to an entire function of exponential type $\le 1$:
\begin{equation}\label{eq:PW}
Y(z)=\int_{-1}^{1}e^{i\xi z}\,d\Phi(\xi),\qquad |Y(z)|\le\|d\Phi\|\,e^{|\Im z|},
\end{equation}
where $d\Phi$ is the Cram\'er spectral measure from \eqref{eq:CramerRep} and $\|d\Phi\|$ is its total variation.
\end{corollary}

\begin{proof}
Define $Y(z):=\int_{-1}^{1}e^{i\xi z}\,d\Phi(\xi)$ for $z\in\mathbb{C}$. The integrand is entire in $z$ for each fixed $\xi\in[-1,1]$; the measure $d\Phi$ is finite (total variation $\|d\Phi\|<\infty$) and the integrand is uniformly bounded on compact subsets of $\mathbb{C}$ by $e^{R}$ for $|z|\le R$ since $|\xi|\le 1$. By the dominated convergence theorem and Morera's theorem, $Y$ is entire. For the exponential-type bound,
\begin{equation}\label{eq:PWbound}
|Y(z)|\le\int_{-1}^{1}|e^{i\xi z}|\,|d\Phi|(\xi)=\int_{-1}^{1}e^{-\xi\Im z}\,|d\Phi|(\xi)\le e^{|\Im z|}\int_{-1}^{1}|d\Phi|(\xi)=\|d\Phi\|\,e^{|\Im z|}.
\end{equation}
Restricted to $z=u\in\mathbb{R}$, \eqref{eq:PW} reproduces \eqref{eq:CramerRep}; uniqueness of the entire extension of the real-analytic function $u\mapsto Y(u)$ (with globally convergent representation) identifies this extension with $Y$ on $\mathbb{R}$.
\end{proof}

\begin{lemma}[Non-negativity of the spectral measure]\label{lem:PSD}
The spectral measure $|d\Phi|^2$ of the representation \eqref{eq:CramerRep} is a non-negative finite Borel measure on $[-1,1]$.
\end{lemma}

\begin{proof}
From \eqref{eq:CramerRep}, the Cram\'er covariance of $Y$ is $C(\eta)=\int_{-1}^{1}e^{i\xi\eta}\,d|\Phi|^2(\xi)$ with $|d\Phi|^2$ non-negative and supported in $[-1,1]$ by construction. The power spectral density $S$ from \eqref{eq:PSDdef} equals the Radon--Nikodym derivative of the absolutely continuous part of $|d\Phi|^2$; Theorem~\ref{thm:bandlim} gives $S(\mu)=0$ for $|\mu|>1$, consistent with the support containment.
\end{proof}

\begin{theorem}[Akhiezer]\label{thm:akhiezer}
A real-valued entire function $F$ of exponential type $\le\tau$ whose spectral measure $|\hat F|^2$ is a non-negative measure supported in $[-\tau,\tau]$ belongs to the Laguerre--P\'olya class. All zeros of $F$ in $\mathbb{C}$ are real.
\end{theorem}

\begin{proof}\cite{Levin}, Lecture 16, Theorem 3; \cite{Akhiezer}, \S V.4.\end{proof}

\begin{corollary}[Reality of zeros of $Y$]\label{cor:Hzero}
All zeros of $Y$ in $\mathbb{C}$ are real.
\end{corollary}

\begin{proof}
$Y$ is real on $\mathbb{R}$ (Definition~\ref{def:Htilde}), entire of exponential type $\le 1$ (Corollary~\ref{cor:PW}), with spectral measure $|d\Phi|^2$ non-negative and supported in $[-1,1]$ (Lemma~\ref{lem:PSD}). Apply Theorem~\ref{thm:akhiezer}.
\end{proof}

\section{Transfer to the Hardy Z-Function on the Real Line}

The transfer uses only the real-line identity
\begin{equation}\label{eq:UTCreal}
Z(t) \;=\; Y(\Theta(t))\,\sqrt{\Theta'(t)},\qquad t\in\mathbb{R},
\end{equation}
which is the rearrangement of \eqref{eq:Ydef} after substituting $u=\Theta(t)$. On $\mathbb{R}$, $\Theta'(t)>0$ (Lemma~\ref{lem:Thetabij}) and $\sqrt{\Theta'(t)}$ is the principal (positive) square root of a strictly positive real number; no branch choice is involved.

\begin{theorem}[Zero-set bijection under the unitary time change on $\mathbb{R}$]\label{thm:UTCzero}
Let $c$ be admissible. The map $t\mapsto\Theta(t)$ is a bijection $\mathbb{R}\to\mathbb{R}$ that restricts to a bijection of zero sets,
\begin{equation}\label{eq:ZeroBij}
\{t\in\mathbb{R}:Z(t)=0\}\;\longleftrightarrow\;\{u\in\mathbb{R}:Y(u)=0\},\qquad u=\Theta(t),
\end{equation}
with preservation of multiplicity: if $t_0\in\mathbb{R}$ and $Y$ has a zero of order $m$ at $\Theta(t_0)$, then $Z$ has a zero of order exactly $m$ at $t_0$, and conversely.
\end{theorem}

\begin{proof}
$\sqrt{\Theta'(t)}$ is strictly positive on $\mathbb{R}$. From \eqref{eq:UTCreal}, for any $t_0\in\mathbb{R}$,
\begin{equation}\label{eq:simpleZero}
Z(t_0)=0 \iff \sqrt{\Theta'(t_0)}\,Y(\Theta(t_0))=0 \iff Y(\Theta(t_0))=0.
\end{equation}
Since $\Theta:\mathbb{R}\to\mathbb{R}$ is a bijection (Lemma~\ref{lem:Thetabij}), $t_0\mapsto\Theta(t_0)$ is a bijection between the real zero set of $Z$ and the real zero set of $Y$.

For multiplicity, let $t_0\in\mathbb{R}$ and suppose $Y$ has a zero of order $m$ at $s_0:=\Theta(t_0)$. Since $Y$ is entire (Corollary~\ref{cor:PW}), there exists an entire $g$ with $g(s_0)\ne 0$ and
\begin{equation}\label{eq:Ysfactor}
Y(s) \;=\; (s-s_0)^m\,g(s).
\end{equation}
Substitute $s=\Theta(t)$ (defined and $C^\infty$ for $t\in\mathbb{R}$):
\begin{equation}\label{eq:YThetafactor}
Y(\Theta(t)) \;=\; (\Theta(t)-\Theta(t_0))^m\,g(\Theta(t)).
\end{equation}
The Hadamard-type identity
\begin{equation}\label{eq:Hadamard}
\Theta(t)-\Theta(t_0) \;=\; (t-t_0)\,h(t),\qquad h(t):=\int_0^1 \Theta'\!\big(t_0+u(t-t_0)\big)\,du,
\end{equation}
valid for all $t\in\mathbb{R}$ because $\Theta\in C^1(\mathbb{R})$, gives $h(t_0)=\Theta'(t_0)>0$, and continuity of $h$ together with $\Theta'>0$ on $\mathbb{R}$ ensures $h(t)>0$ for $t$ in a neighborhood of $t_0$. Combining \eqref{eq:UTCreal}, \eqref{eq:YThetafactor}, and \eqref{eq:Hadamard},
\begin{equation}\label{eq:Zfactor}
Z(t) \;=\; Y(\Theta(t))\sqrt{\Theta'(t)} \;=\; (t-t_0)^m\,h(t)^m\,g(\Theta(t))\,\sqrt{\Theta'(t)}.
\end{equation}
Denote the trailing factor
\begin{equation}\label{eq:Hfactor}
H(t) \;:=\; h(t)^m\,g(\Theta(t))\,\sqrt{\Theta'(t)}.
\end{equation}
Then $H$ is continuous at $t_0$ and
\begin{equation}\label{eq:Ht0}
H(t_0) \;=\; \Theta'(t_0)^m\,g(s_0)\,\sqrt{\Theta'(t_0)} \;\neq\; 0
\end{equation}
because $\Theta'(t_0)>0$ and $g(s_0)\neq 0$. Therefore $Z$ has a zero of order exactly $m$ at $t_0$.

The converse (given a zero of $Z$ of order $m$ at $t_0\in\mathbb{R}$, $Y$ has a zero of order $m$ at $\Theta(t_0)$) is obtained by applying the same argument with $\Theta$ replaced by $\Theta^{-1}$, which satisfies $(\Theta^{-1})'(u)=1/\Theta'(\Theta^{-1}(u))>0$ for $u\in\mathbb{R}$ by Lemma~\ref{lem:Thetabij} and the inverse function theorem.
\end{proof}

\begin{corollary}[Zeros of $Z$ on $\mathbb{R}$]\label{cor:Zreal}
The zero set of $Z$ on $\mathbb{R}$, counted with multiplicity, is $\Theta^{-1}$ applied to the zero set of $Y$, counted with multiplicity.
\end{corollary}

\begin{proof}
Corollary~\ref{cor:Hzero} gives that every zero of the entire function $Y$ lies on $\mathbb{R}$. Theorem~\ref{thm:UTCzero} gives a multiplicity-preserving bijection between the zero set of $Y$ on $\mathbb{R}$ and the zero set of $Z$ on $\mathbb{R}$, implemented by $t\mapsto\Theta(t)$.
\end{proof}

\section{The Zero Locus of $\zeta$ on the Critical Line}

\begin{corollary}[Zeros of $\zeta$ on the critical line]\label{cor:Zeta-crit}
The zero set of $\zeta$ on the critical line $\{\Re s=\tfrac{1}{2}\}$, counted with multiplicity, is the image of $\Theta^{-1}(\{u\in\mathbb{R}:Y(u)=0\})$ under the map $t\mapsto\tfrac{1}{2}+it$, with multiplicities inherited from those of $Y$.
\end{corollary}

\begin{proof}
For $t\in\mathbb{R}$, $|e^{i\theta(t)}|=1$, so $e^{i\theta(t)}$ is a nonzero complex number; hence $\zeta(\tfrac{1}{2}+it)=e^{-i\theta(t)}Z(t)$ gives $\zeta(\tfrac{1}{2}+it)=0 \iff Z(t)=0$, with matching orders of vanishing at every $t\in\mathbb{R}$ (multiplication by the non-vanishing real-analytic factor $e^{-i\theta(t)}$ does not alter the order of a zero of the real-analytic factor $Z(t)$ at a real point). Corollary~\ref{cor:Zreal} identifies the real zero set of $Z$ with $\Theta^{-1}$ applied to the zero set of $Y$. The map $t\mapsto\tfrac{1}{2}+it$ is a bijection $\mathbb{R}\to\{\Re s=\tfrac{1}{2}\}$. Composing gives the stated identification.
\end{proof}

\begin{corollary}[Spectral reconstruction of $\zeta$ on the critical line]\label{cor:Zeta-recon}
For every $t\in\mathbb R$,
\begin{equation}\label{eq:Zeta-recon}
\zeta(\tfrac{1}{2}+it)\;=\;e^{-i\theta(t)}\,\sqrt{\Theta'(t)}\,\int_{-1}^{1}e^{i\xi\,\Theta(t)}\,d\Phi(\xi),
\end{equation}
where $d\Phi$ is the spectral measure from Corollary~\ref{cor:cramer} with the per-mode factorization \eqref{eq:Phifactor}.
\end{corollary}

\begin{proof}
$\zeta(\tfrac{1}{2}+it)=e^{-i\theta(t)}Z(t)$ on $\mathbb R$ (Introduction). $Z(t)=Y(\Theta(t))\sqrt{\Theta'(t)}$ on $\mathbb R$ by \eqref{eq:UTCreal}. $Y(u)=\int_{-1}^{1}e^{i\xi u}\,d\Phi(\xi)$ for $u\in\mathbb R$ by \eqref{eq:CramerRep}. Compose with $u=\Theta(t)$.
\end{proof}

\appendix
\begin{thebibliography}{9}

\bibitem{AbramowitzStegun}
M.~Abramowitz and I.~A.~Stegun (eds.), \emph{Handbook of Mathematical Functions}, Dover, 1972.

\bibitem{ClercMallat2003}
M.~Clerc and S.~Mallat, \emph{Estimating deformations of stationary processes}, The Annals of Statistics \textbf{31} (2003), no.~6, 1772--1821.

\bibitem{CramerLeadbetter}
H.~Cram\'er and M.~R.~Leadbetter, \emph{Stationary and Related Processes}, Wiley, 1967.

\bibitem{Akhiezer}
N.~I.~Akhiezer, \emph{Theory of Approximation}, Ungar, 1956.

\bibitem{Edwards}
H.~M.~Edwards, \emph{Riemann's Zeta Function}, Academic Press, 1974.

\bibitem{Hormander}
L.~H\"ormander, \emph{The Analysis of Linear Partial Differential Operators I}, 2nd ed., Springer, 1990.

\bibitem{Levin}
B.~Ya.~Levin, \emph{Lectures on Entire Functions}, Translations of Mathematical Monographs, vol.~150, American Mathematical Society, 1996.

\bibitem{Rudin}
W.~Rudin, \emph{Real and Complex Analysis}, 3rd ed., McGraw--Hill, 1987.

\bibitem{Rudin-FA}
W.~Rudin, \emph{Functional Analysis}, 2nd ed., McGraw--Hill, 1991.

\bibitem{Titchmarsh}
E.~C.~Titchmarsh, \emph{The Theory of the Riemann Zeta-Function}, 2nd ed., revised by D.~R.~Heath-Brown, Oxford University Press, 1986.

\end{thebibliography}

\end{document}
